% Options for packages loaded elsewhere
\PassOptionsToPackage{unicode}{hyperref}
\PassOptionsToPackage{hyphens}{url}
\PassOptionsToPackage{dvipsnames,svgnames,x11names}{xcolor}
%
\documentclass[
]{letter}

\usepackage{amsmath,amssymb}
\usepackage{iftex}
\ifPDFTeX
  \usepackage[T1]{fontenc}
  \usepackage[utf8]{inputenc}
  \usepackage{textcomp} % provide euro and other symbols
\else % if luatex or xetex
  \usepackage{unicode-math}
  \defaultfontfeatures{Scale=MatchLowercase}
  \defaultfontfeatures[\rmfamily]{Ligatures=TeX,Scale=1}
\fi
\usepackage{lmodern}
\ifPDFTeX\else  
    % xetex/luatex font selection
\fi
% Use upquote if available, for straight quotes in verbatim environments
\IfFileExists{upquote.sty}{\usepackage{upquote}}{}
\IfFileExists{microtype.sty}{% use microtype if available
  \usepackage[]{microtype}
  \UseMicrotypeSet[protrusion]{basicmath} % disable protrusion for tt fonts
}{}
\makeatletter
\@ifundefined{KOMAClassName}{% if non-KOMA class
  \IfFileExists{parskip.sty}{%
    \usepackage{parskip}
  }{% else
    \setlength{\parindent}{0pt}
    \setlength{\parskip}{6pt plus 2pt minus 1pt}}
}{% if KOMA class
  \KOMAoptions{parskip=half}}
\makeatother
\usepackage{xcolor}
\usepackage[margin=1in,bottom=1in,left=1in,right=1in]{geometry}
\setlength{\emergencystretch}{3em} % prevent overfull lines
\setcounter{secnumdepth}{-\maxdimen} % remove section numbering
% Make \paragraph and \subparagraph free-standing
\makeatletter
\ifx\paragraph\undefined\else
  \let\oldparagraph\paragraph
  \renewcommand{\paragraph}{
    \@ifstar
      \xxxParagraphStar
      \xxxParagraphNoStar
  }
  \newcommand{\xxxParagraphStar}[1]{\oldparagraph*{#1}\mbox{}}
  \newcommand{\xxxParagraphNoStar}[1]{\oldparagraph{#1}\mbox{}}
\fi
\ifx\subparagraph\undefined\else
  \let\oldsubparagraph\subparagraph
  \renewcommand{\subparagraph}{
    \@ifstar
      \xxxSubParagraphStar
      \xxxSubParagraphNoStar
  }
  \newcommand{\xxxSubParagraphStar}[1]{\oldsubparagraph*{#1}\mbox{}}
  \newcommand{\xxxSubParagraphNoStar}[1]{\oldsubparagraph{#1}\mbox{}}
\fi
\makeatother


\providecommand{\tightlist}{%
  \setlength{\itemsep}{0pt}\setlength{\parskip}{0pt}}\usepackage{longtable,booktabs,array}
\usepackage{calc} % for calculating minipage widths
% Correct order of tables after \paragraph or \subparagraph
\usepackage{etoolbox}
\makeatletter
\patchcmd\longtable{\par}{\if@noskipsec\mbox{}\fi\par}{}{}
\makeatother
% Allow footnotes in longtable head/foot
\IfFileExists{footnotehyper.sty}{\usepackage{footnotehyper}}{\usepackage{footnote}}
\makesavenoteenv{longtable}
\usepackage{graphicx}
\makeatletter
\newsavebox\pandoc@box
\newcommand*\pandocbounded[1]{% scales image to fit in text height/width
  \sbox\pandoc@box{#1}%
  \Gscale@div\@tempa{\textheight}{\dimexpr\ht\pandoc@box+\dp\pandoc@box\relax}%
  \Gscale@div\@tempb{\linewidth}{\wd\pandoc@box}%
  \ifdim\@tempb\p@<\@tempa\p@\let\@tempa\@tempb\fi% select the smaller of both
  \ifdim\@tempa\p@<\p@\scalebox{\@tempa}{\usebox\pandoc@box}%
  \else\usebox{\pandoc@box}%
  \fi%
}
% Set default figure placement to htbp
\def\fps@figure{htbp}
\makeatother

\usepackage{booktabs}
\usepackage{caption}
\usepackage{longtable}
\usepackage{colortbl}
\usepackage{array}
\usepackage{anyfontsize}
\usepackage{multirow}
\makeatletter
\@ifpackageloaded{caption}{}{\usepackage{caption}}
\AtBeginDocument{%
\ifdefined\contentsname
  \renewcommand*\contentsname{Table des matières}
\else
  \newcommand\contentsname{Table des matières}
\fi
\ifdefined\listfigurename
  \renewcommand*\listfigurename{Liste des Figures}
\else
  \newcommand\listfigurename{Liste des Figures}
\fi
\ifdefined\listtablename
  \renewcommand*\listtablename{Liste des Tables}
\else
  \newcommand\listtablename{Liste des Tables}
\fi
\ifdefined\figurename
  \renewcommand*\figurename{Figure}
\else
  \newcommand\figurename{Figure}
\fi
\ifdefined\tablename
  \renewcommand*\tablename{Table}
\else
  \newcommand\tablename{Table}
\fi
}
\@ifpackageloaded{float}{}{\usepackage{float}}
\floatstyle{ruled}
\@ifundefined{c@chapter}{\newfloat{codelisting}{h}{lop}}{\newfloat{codelisting}{h}{lop}[chapter]}
\floatname{codelisting}{Listing}
\newcommand*\listoflistings{\listof{codelisting}{Liste des Listings}}
\makeatother
\makeatletter
\makeatother
\makeatletter
\@ifpackageloaded{caption}{}{\usepackage{caption}}
\@ifpackageloaded{subcaption}{}{\usepackage{subcaption}}
\makeatother

\ifLuaTeX
\usepackage[bidi=basic]{babel}
\else
\usepackage[bidi=default]{babel}
\fi
\babelprovide[main,import]{french}
% get rid of language-specific shorthands (see #6817):
\let\LanguageShortHands\languageshorthands
\def\languageshorthands#1{}
\usepackage{bookmark}

\IfFileExists{xurl.sty}{\usepackage{xurl}}{} % add URL line breaks if available
\urlstyle{same} % disable monospaced font for URLs
\hypersetup{
  pdfauthor={Pierre J. FISCHER},
  pdflang={fr},
  pdfsubject={Assignation en partage judiciaire de succession FISCHER
née DURRENBERGER Aline},
  colorlinks=true,
  linkcolor={blue},
  filecolor={Maroon},
  citecolor={Blue},
  urlcolor={Blue},
  pdfcreator={LaTeX via pandoc}}


\author{Pierre J. FISCHER}
\date{24 juillet 2026}

\begin{document}
\signature{Pierre J. FISCHER}
\address{Pierre J. FISCHER\\fishrp@pm.me +33 76748 2009\\39 rue de la
Figairasse, Bat I\\MONTPELLIER 34070\\F-}
\begin{letter}{Tribunal de Proximité de Haguenau\\41 rue de la
Redoute\\CS 10240\\67504 HAGUENAU Cedex\\ ~ \\Objet: Assignation en
partage judiciaire de succession FISCHER née DURRENBERGER Aline}
\opening{Madame, Monsieur,}


LE CONTEXTE

Je soussigné, Pierre Jacques FISCHER, demeurant au 39 rue de la
Figairasse, Bât I, 34070 MONTPELLIER, ai l'honneur de vous saisir par la
présente afin de solliciter un partage judiciaire concernant la
succession de feu FISCHER née DURRENBERGER Aline, décédée selon
l'\textbf{acte de décès}\(^d\) de FISCHER née DURRENBERGER Aline du 19
Nov 2021, Numéro 37 - Mairie de Mertzwiller

En effet, trois années après le décès et deux tentatives de règlement
amiable, il n'a pas été possible de parvenir à un accord entre les
parties concernées. Par conséquent, je me vois contraint de demander
l'intervention du tribunal pour procéder au partage des biens
conformément aux dispositions légales en vigueur.

En effet, les trois soeurs ont refusé tout à la fois, de faire la
réunion fictive des biens distribués après la donation-partage de 1986
et de me restituer, sur l'immeuble sis à Mulhouse, le quasi-usufruit né
d'un partage judiciaire qui m'a été imposé en 2011. Et alors que le
notaire refuse de m'adresser un procès-verbal de carence, la situation
est dans l'impasse.

LES PARTIES

la \textbf{donation-partage}\(^a\) du 19 Fev 1986 des époux Georges
Fischer et Aline Durrenberger, par Me Lotz en quatre parts égales de
biens immobiliers au profit des parties héritières et seuls quatre
enfants, identifiés par le projet \textbf{d'affirmation
sacramentelle}\(^b\) du 3 mars 2022 qui actualise les adresses comme
suit:

\begin{itemize}
\item
  Michèle Fischer, demeurant à Mertzwiller (67580) 9 rue du Général
  Koenig
\item
  Marguerite Besnier, demeurant à Goersdorf (67360) 9 rue des
  châtaigniers
\item
  Aline Salzeman, demeurant à Armeau (89500) 34 rue de l'île de France
\item
  Pierre Fischer, demeurant à Montpellier (34070) 39 rue de la
  Figairasse Bât I

  ~\\
\end{itemize}

LE PATRIMOINE EXISTANT

\begin{itemize}
\item
  en plus de la donation-partage sus-mentionnée la \textbf{déclaration
  de succession}\(^e\) du 3 mars 2022 projetée par Me Lotz, se limitait
  au partage en quatre parts égales
\item
  du seul \textbf{actif net successoral} à hauteur de \textbf{129 026,58
  €}. Cette proposition n'a pas trouvé l'accord de Pierre F. pour les
  raisons ci-dessous.
\end{itemize}

LES MOTIFS

\begin{itemize}
\item
  la demande de restitution de \textbf{74700 €} correspondant à
  \textbf{l'usufruit perdu}\(^c\) précocemment par la vente judiciaire
  forcée en 2012 du lot immobilier de Pierre F. (le nu-proprétaire) se
  heurte à un refus des trois soeurs alors même qu'elles défendent
  Michèle F (nue-propriétaire) qui a \textbf{consommé l'usufruit} sur
  son lot avant le décès du \emph{de cujus.} Considérant l'égalité dans
  la nature des biens (i.e.~l'usufruit sur le bien immobiliier dans la
  période entre la donation et le décès ), ainsi que l'égalité ès
  qualité d'héritier, il n'est pas recevable de leur appliquer \emph{in
  fine} une règle de partage doublement discriminatoire et inique.
\item
  la lettre \textbf{lotz\_rapport\_signe.pdf}\(^g\) du 21 Juin 2024
  demande le \textbf{rapport des libéralités,} sérieusement étayé par le
  chiffrage et l'argumentaire du document
  \textbf{lotz\_technique-v0.99.pdf}\(^h\) qui analyse les factures
  afférentes aux gros chèques (préalablment identifiés sur les relevés
  bancaires), et permet d'identifier la nature, le montant et les
  bénéficiaires des gros travaux ainsi que le sort de l'usufruit. Cette
  étude conclut à un transfert massif et tardif de \textbf{libéralités
  déguisées} aux profit des seules soeurs concomittant à
  l'appauvrissement de la donatrice.
\end{itemize}

\emph{en résumé, \textbf{le rapport des libéralités} (dont l'effet
touchera l'ensemble des héritiers) demande à retenir ce qui est
classique en pareil cas, à savoir les investissements en gros travaux
(133 K€), le bénéfice du logement gratuit sur 35 ans (160 K€), la valeur
du mobilier restant (30 K€) et la maison (160 K€) érigée après la
donation de 1986 sur le terrain nu sis à Mertzwiller 9 rue d'Eschbach
(lot de Marguerite B.) et de refaire en conséquence le partage de la
réunion fictive ainsi obtenue. Cf détail sous ``Demandes''.}

\begin{itemize}
\tightlist
\item
  enfin, \textbf{l'argument du faisceau} réside dans l'observation que
  les avantages (dont la matérialité non négigeable se monte à 450
  K\emph{€)} sont toujours au bénéfice des soeurs et au détriment du
  fils.
\end{itemize}

LES DILIGENCES

\begin{itemize}
\tightlist
\item
  la \textbf{lettre\_rappel\_signe.pdf}\(^f\) résume l'histoire des
  diligences, restées vaines du 15 nov 2023 et 22 juin 2024 alors que
\end{itemize}

\emph{le dernier échange de juin 2024, n'a reçu de Me Lotz qu'une fin de
non-recevoir, actée en retour par ma lettre
\textbf{2024-10-11\_declinaison\_signée.pdf}}\(^i\) \emph{qui met fin à
trois années d'effort amiable.}

\begin{itemize}
\tightlist
\item
  enfin \textbf{LETRECO\_Depot\_R5OK2POF\_Successio-enberger.pdf}\(^j\)
  du 22 Juin 2024 apporte la preuve légale de cette diligence par la
  transmission à Me Lotz de \textbf{lotz\_technique-v0.99.pdf}
  détaillant la source, les méthodes d'évaluation ainsi que les
  libéralités qualifiables de donation déguisées, que ma demande dans
  \textbf{lotz\_rapport\_signe.pdf} souhaite voir rapporter à ce projet
  de partage revisé.
\end{itemize}

LES DEMANDES

Si la donation-partage de 1986 était voulue égalitaire en son temps, je
soussigné Pierre F. demande la réévaluation de la situation au moment du
décès par la double prise en considération a) du
\textbf{``quasi-usufruit'' perdu} et b) du \textbf{rapport des
libéralités} dûment étayé. Pour ces raisons, je demande au tribunal

\begin{enumerate}
\def\labelenumi{\arabic{enumi}.}
\tightlist
\item
  en première instance d'accorder la \textbf{restitution du
  quasi-usufruit} au profit de Pierre F. pour une valeur de 74.7 K€ au
  motif de lever le caractère discriminatoire dans le traitement des
  usufruits consommés avant le décès.
\end{enumerate}

\begin{quote}
\emph{En effet il n'est pas recevable de faire valoir pour s'y opposer
d'invoquer une double iniquité: celle de refuser la restitution d'un
usufruit perdu par Pierre F et de défendre en même temps l'usufruit
gagné par Michèle F.}
\end{quote}

\begin{enumerate}
\def\labelenumi{\arabic{enumi}.}
\setcounter{enumi}{1}
\tightlist
\item
  d'effectuer ensuite, la \textbf{réunion fictive} de l'actif restant
  avec le rapport des libéralités chiffrées par le document
  lotz\_technique et d'opérer le partage sur le modèle ci-dessous (ici
  sans la restitution préalalble de l'usufruit).
\end{enumerate}

\begin{table}
\caption*{
{\fontsize{20}{25}\selectfont  Formule de partage amiable\fontsize{12}{15}\selectfont }
} 
\fontsize{12.0pt}{14.0pt}\selectfont
\begin{tabular*}{\linewidth}{@{\extracolsep{\fill}}l|lrrr}
\toprule
 & nature & deja\_recu & quotepart & a\_recevoir \\ 
\midrule\addlinespace[2.5pt]
\multicolumn{5}{l}{michèle} \\[2.5pt] 
\midrule\addlinespace[2.5pt]
sous\_total & — & -191.50 & 145.21 & -46.29 \\ 
\midrule 
 & loyer & -160.0 & 48.4 & -111.6 \\ 
 & mobilier & -30.0 & 48.4 & 18.4 \\ 
 & travaux & -1.5 & 48.4 & 46.9 \\ 
\midrule\addlinespace[2.5pt]
\multicolumn{5}{l}{marguerite} \\[2.5pt] 
\midrule\addlinespace[2.5pt]
sous\_total & — & -187.00 & 145.21 & -41.79 \\ 
\midrule 
 & travaux & -27.0 & 72.6 & 45.6 \\ 
 & maison & -160.0 & 72.6 & -87.4 \\ 
\midrule\addlinespace[2.5pt]
\multicolumn{5}{l}{aline} \\[2.5pt] 
\midrule\addlinespace[2.5pt]
sous\_total & — & -71.00 & 145.21 & 74.21 \\ 
\midrule 
 & travaux & -71.0 & 145.2 & 74.2 \\ 
\midrule\addlinespace[2.5pt]
\multicolumn{5}{l}{pierre} \\[2.5pt] 
\midrule\addlinespace[2.5pt]
sous\_total & — & -2.34 & 145.21 & 142.87 \\ 
\midrule 
 & don & -2.3 & 145.2 & 142.9 \\ 
\midrule 
\midrule 
reunion & — & -451.84 & 580.84 & 129 \\ 
\bottomrule
\end{tabular*}
\end{table}

\begin{quote}
\emph{De surcroît, les libéralités supplémentaires accordées aux soeurs
après la perte de l'usufruit de Pierre F. en 2012, n'ont fait que
creuser cette inégalité, toujours à mon détriment, sans aucune vélleité
de compensation de la part de Fischer Aline.}
\end{quote}

\begin{enumerate}
\def\labelenumi{\arabic{enumi}.}
\setcounter{enumi}{2}
\tightlist
\item
  le remboursement des \textbf{dépens} ainsi que des \textbf{frais
  irrépétibles} selon CPC Art 700.
\end{enumerate}

LE CHOIX DU NOTAIRE

Je souhaite voir commis le
\href{https://www.bing.com/ck/a?!&&p=f2066ee7e9a9fe535e713659e0db4ee90bd39e3fa6431940c2eba6275c4f6b87JmltdHM9MTczOTE0NTYwMA&ptn=3&ver=2&hsh=4&fclid=04d29d0e-2a39-6ad9-1b33-89a22b806b7b&psq=notaire+Schott+strabsourg&u=a1aHR0cHM6Ly9zY2hvdHQtc2Nod2FhYi1naWxsZXQtbm90YWlyZXMuZnIv&ntb=1}{Notaire
Valentin Schott ou à défault Yannick Schott} 1 Rue du Dôme - 67000
Strasbourg. etude@67006.notaires.fr \href{tel:+33388321087}{03 88 32 10
87}

Restant à votre disposition pour toute information complémentaire.



\closing{Je vous prie d'agréer, Madame, Monsieur, l'expression de mes
salutations distinguées,}
\vfill
\encl{19 Fev 1986 a\_2\_DONATION\_FISCHER\_1986.pdf\\3 Mars 2022
b\_PROJET\_AFFIRMATION SACRAMENTELLE FISCHER Aline née
DURRENBERGER.pdf\\25 Sep 2012 c\_D\_decompte\_mulhouse.pdf
(usufruit)\\19 Nov 2021 d\_1\_DECES\_epoux\_FISCHER.pdf\\3 Mars 2022
e\_PROJET\_DECLARATION SUCCESSION FISCHER Aline née DURRENBERGER.pdf\\21
Jun 2024 f\_lettre\_rappel\_signe.pdf\\21 Jun 2024
g\_lotz\_rapport\_signe.pdf\\21 Jun 2024 h\_lotz\_technique-v0.99.pdf
(22 pages)\\11 Oct 2024 i\_2024-10-11\_declinaison\_signée.pdf\\22 Jun
2024 j\_LETRECO\_Depot\_R5OK2POF\_Successio-enberger.pdf\\10 Fev 2025
k\_saisine\_haguenau (la présente)}
\end{letter}
\end{document}
